\chapter{Game Solutions and Matching as Argumentation}
\label{sec:p18-game-argumentation}

The three correspondences in Dung's Section~3 concern
coalitional domination between payoff vectors, the
undominated core, and blocking pairs in a matching
\cite[pp.~334--338]{Dung1995}. Their solution objects
and deviation rules are different. We state their
feasibility and preference assumptions before
constructing their attack graphs.

\section{Cooperative games and sets of imputations}

\begin{hbdefinition}{def:p18-cooperative-game}{Transferable-utility games and feasible imputations}
Let \(N=\{1,\ldots,n\}\), \(n\geq1\), be the
finite player set. A characteristic-function
game assigns a real value \(v(K)\) to each
coalition \(K\subseteq N\), with
\(v(\varnothing)=0\). Assume superadditivity:
\[
       v(K\cup L)\geq v(K)+v(L)\qquad(K\cap L=\varnothing).
\]
The values and payoffs use a common transferable
unit. They are not arbitrary ordinal ranks.
For \(x\in\mathbb R^n\), write
\(x(K)=\sum_{i\in K}x_i\). Its imputation set is
\[
 I(v)=\{x\in\mathbb R^n:x(N)=v(N),\
                                  x_i\geq v(\{i\})\text{ for every }i\}.
\]
The two conditions are efficiency and individual
rationality. The complete primary definition
appears in Lucas's treatment
\cite[sec.~2, pp.~1--2]{Lucas1968NoSolutionProof}.

For \(x,y\in I(v)\), write \(x\succ_K y\)
when \(K\ne\varnothing\),
\[
            x_i>y_i\quad(i\in K),\qquad x(K)\leq v(K).
\]
The strict improvement concerns each member
of the coalition; the feasibility inequality
is weak. Write \(x\succ_v y\) if some
coalition witnesses this domination.
\end{hbdefinition}

\begin{hblemma}{thm:p18-imputation-core}{proved=annex:p18.imputation.core}{Feasibility and the two equivalent core definitions}
The imputation set is nonempty. Domination
is irreflexive. Its undominated core
\[
 {\cal C}(v)=\{y\in I(v):\text{no }x\in I(v)\text{ dominates }y\}
\]
is exactly
\[
              \{y\in I(v):y(K)\geq v(K)\text{ for every }K\subseteq N\}.
\]
The proof explicitly extends any feasible
blocking coalition to an individually
rational efficient imputation.
\end{hblemma}

\begin{hbdefinition}{def:p18-nm-solutions}{Internal and external stability}
A von Neumann--Morgenstern solution, also
called an NM solution or stable set, is
a set \(S\subseteq I(v)\) satisfying
\[
\begin{aligned}
 &\text{no }x,y\in S\text{ have }x\succ_v y,\\
 &\text{for every }y\in I(v)\setminus S
                       \text{ some }x\in S\text{ has }x\succ_v y.
\end{aligned}
\]
The first condition is internal stability
and the second external stability.
Construct the framework
\[
             {\cal A}_v=(I(v),\{(x,y):x\succ_v y\}).
\]
A preferred solution in Dung's additional
sense means a preferred extension of this
framework. It is a maximal admissible
set of imputations, not a payoff vector
of maximum utility.
\end{hbdefinition}

\begin{hbtheorem}{thm:p18-nm-argumentation}{proved=annex:p18.nm.argumentation}{Both cooperative-game correspondences}
The NM solutions of the game are exactly
the stable extensions of \({\cal A}_v\).
If \(F_v\) is its characteristic operator,
\[
                       {\cal C}(v)=F_v(\varnothing).
\]
Consequently the core is contained in the
grounded extension and in every NM solution.
Preferred solutions exist by arbitrary-framework
preferred existence. This does not establish
existence of an NM solution. The imputation
set may be uncountable although the player
set is finite.
\end{hbtheorem}

\begin{hbdefinition}{def:p18-three-player-game}{The source's three-player game}
Take \(N=\{1,2,3\}\), with
\[
 v(\varnothing)=v(\{i\})=0,\qquad
                 v(K)=2\quad(|K|\geq2).
\]
Then \(I(v)=\{x:x_i\geq0,\ x_1+x_2+x_3=2\}\).
Define
\[
 S_\triangle=\{(1,1,0),(1,0,1),(0,1,1)\},
 \qquad
 S_c=\{(t,2-c-t,c):0\leq t\leq2-c\}\quad(0\leq c\leq2).
\]
The fixed-share family makes the source's
second solution description explicit.
For \(0\leq c<2/3\), player 3 receives
less than the equal-share benchmark \(2/3\)
at every imputation in that solution.
\end{hbdefinition}

\begin{hbtheorem}{thm:p18-three-player-solutions}{proved=annex:p18.three.player.solutions}{Both source solution families and every endpoint case}
The displayed game is superadditive.
\(S_\triangle\) is an NM solution.
For \(0\leq c\leq2\), \(S_c\) is an
NM solution exactly when \(c<1\).
Thus the game has infinitely many
distinct NM solutions, including the
source's fixed-share cases \(c<2/3\).
Its core and grounded extension are
empty. Domination in this game is
not transitive.
\end{hbtheorem}

\section{Strict two-sided matching}

\begin{hbdefinition}{def:p18-marriage-model}{The exact balanced complete-list matching model}
Let \(M,W\) be disjoint finite sets with
\(|M|=|W|=n\), allowing \(n=0\).
Each \(m\in M\) has a strict total
order \(\succ_m\) on all of \(W\),
and each \(w\in W\) has a strict
total order \(\succ_w\) on all of \(M\).
Thus each order is irreflexive and transitive,
and compares every two distinct partners.
There are no ties or unacceptable partners.
A matching is a bijection \(\mu:M\to W\),
identified with its graph in \(M\times W\).
It is stable when no unchosen pair
\((m,w)\) satisfies
\[
                        w\succ_m\mu(m),\qquad
                        m\succ_w\mu^{-1}(w).
\]
Such a pair would be a blocking pair.

The argument set is \(M\times W\), with
\[
 (m,w)\to(m',w')\quad\Longleftrightarrow\quad
 \bigl(m=m'\land w\succ_m w'\bigr)
       \lor\bigl(w=w'\land m\succ_w m'\bigr).
\]
The attacker is the shared participant's
preferred proposed pairing.
\end{hbdefinition}

\begin{hbtheorem}{thm:p18-marriage-argumentation}{proved=annex:p18.marriage.argumentation}{The complete matching/stable-extension equivalence}
A set of these pair arguments is
conflict-free exactly when it is
a partial matching. Its stable
extensions are exactly the graphs
of stable bijections \(M\to W\).
The reverse direction proves
totality using finiteness and
equal cardinalities. The empty
instance has the unique empty
matching and the unique stable
extension \(\varnothing\).
\end{hbtheorem}

\begin{hbmethod}{def:p18-deferred-acceptance}{Deferred acceptance with proposals from \(M\)}
Initially nobody is held. Repeatedly choose
an unheld \(m\) who has an unproposed
partner. He proposes to his highest-ranked
unproposed \(w\). She keeps the better
of her current held proposer, if any,
and the new proposer, rejecting the
other. A rejection frees that proposer.
Each pair is proposed at most once.
When there is no unheld proposer,
the held pairs are the output.
\end{hbmethod}

\begin{hbtheorem}{thm:p18-deferred-acceptance}{proved=annex:p18.deferred.acceptance}{Termination, existence and the exact optimality claim}
On the stated finite balanced complete-list
domain the procedure terminates after
at most \(n^2\) proposals and returns
a stable matching. Each proposer receives
his most-preferred partner among all
stable matchings. Every receiver receives
her least-preferred partner among all
stable matchings.

The proposer-optimal matching is unique,
so the choice of which free proposer
acts next does not change the output.
Reversing the roles gives the receiver-optimal
matching. These are side-specific
optimality statements, not maximization
of aggregate welfare
\cite[secs.~3 and 5]{GaleShapley1962}.
\end{hbtheorem}

\begin{hbdefinition}{def:p18-matching-boundary-examples}{Three changes to the matching assumptions}
\begin{enumerate}
\item With two participants on each side and all
      alternatives tied, use the same strict-attack
      formula, which gives no attacks.
\item With one participant on each side and no
      mutually acceptable pair, use only acceptable
      pairs as arguments, giving the empty graph.
\item Let \(M=\{m_0,m_1,\ldots\}\) and
      \(W=\{w_0,w_1,\ldots\}\). Each \(m_{i+1}\)
      ranks \(w_i\) first, and each \(w_i\)
      ranks \(m_{i+1}\) first. Complete the
      remaining rankings in increasing index
      order; \(m_0\) ranks all women in
      increasing index order.
\end{enumerate}
\end{hbdefinition}

\begin{hbtheorem}{thm:p18-matching-assumption-boundaries}{proved=annex:p18.matching.assumption.boundaries}{Why the exact matching hypotheses matter}
In the tied example the entire four-pair
argument set is stable but is not a
matching. In the unacceptable-pair
example the empty set is stable but
there is no acceptable perfect matching.
In the infinite example
\[
                         S=\{(m_{i+1},w_i):i<\omega\}
\]
is a stable extension of the strict
complete-list graph, but leaves \(m_0\)
unmatched. Thus equal infinite cardinalities
do not replace the finite totality
argument.
\end{hbtheorem}

\section{One-population pairing and the five-person example}

\begin{hbdefinition}{def:p18-roommate-model}{Distinct unordered pairs and partial matching stability}
Let \(P\) be a finite set. Each person
strictly ranks every other person, and
regards every such partner as preferable
to being unmatched, denoted \(\bot\).
Use the pair set
\[
                    {P\choose2}=\{\{a,b\}:a,b\in P,\ a\ne b\}.
\]
Two distinct pair arguments sharing a
person \(a\) attack according to that
person's preference:
\[
              \{a,b\}\to\{a,c\}\quad\Longleftrightarrow\quad b\succ_a c.
\]
A partial matching is a set of disjoint
pairs. For it, \(\mu(a)\) is the
partner of \(a\), or \(\bot\) if
unmatched. It is stable when no
unchosen distinct pair \(\{a,b\}\)
has both \(b\succ_a\mu(a)\) and
\(a\succ_b\mu(b)\).
\end{hbdefinition}

\begin{hbtheorem}{thm:p18-roommate-argumentation}{proved=annex:p18.roommate.argumentation}{The exact one-population correspondence}
Conflict-free sets are exactly partial
matchings. Stable extensions are exactly
stable partial matchings in the stated
sense. A stable partial matching leaves
at most one person unmatched.
Preferred extensions are maximal admissible
sets of pair arguments; they need not
maximize the number of people matched.
\end{hbtheorem}

\begin{hbdefinition}{def:p18-five-person-pairing}{The source's five participants with their full preference scope}
Let \(P=\{m,w,p_1,p_2,p_3\}\), all distinct.
Require \(m,w\) to rank each other first,
and require the first choices
\[
                     p_1:p_2,\qquad p_2:p_3,\qquad p_3:p_1.
\]
Every remaining ranking is any strict
completion over the other four persons.
Use the unordered-pair graph just defined.

For comparison define an ordered-role
graph on \(P\times P\), giving each
person herself or himself last place
to make diagonal comparisons defined,
and using only the source's same-first
or same-second-coordinate attack rule.
These ordered arguments represent two
roles per person, and are not identified
with unordered marriages.
\end{hbdefinition}

\begin{hbtheorem}{thm:p18-five-person-pairing}{proved=annex:p18.five.person.pairing}{The corrected example and the literal ordered-role boundary}
For every stated strict completion, the
unordered-pair graph has exactly the
admissible sets
\[
                            \varnothing,\qquad \{\{m,w\}\}.
\]
Its grounded and unique preferred
extension is \(\{\{m,w\}\}\), and
it has no stable extension.
The maximum cardinality of an ordinary
partial matching is two pairs, so
the preferred set does not attain it.

In the ordered-role graph, however,
\[
 \{(m,w),(w,m),(p_1,p_2),(p_2,p_3),(p_3,p_1)\}
\]
is a stable extension. It remains
stable if diagonal pairs are deleted.
Thus the unordered construction is
essential for the claimed physical
pairing example; reusing \(P\times P\)
unchanged gives a different result.
\end{hbtheorem}

\begin{hbdefinition}{def:p18-matching-welfare-example}{Two stable matchings and explicit numerical scales}
Let \(M=\{a,b\}\), \(W=\{x,y\}\), with
\[
 a:x\succ y,\quad b:y\succ x,\qquad
 x:b\succ a,\quad y:a\succ b.
\]
Give each proposer utility 1 for his
first choice and 0 for his second;
give each receiver utility 10 for her
first choice and 0 for her second.
These numbers are additional stipulated
scales representing the stated orders.
Write \(U(\eta)\) for their sum under
a matching \(\eta\). A matching Pareto-dominates
another when every participant weakly prefers
it and at least one strictly prefers it.
\end{hbdefinition}

\begin{hbtheorem}{thm:p18-matching-welfare-boundary}{proved=annex:p18.matching.welfare.boundary}{Stability, side-optimality and aggregate welfare differ}
Both perfect matchings
\[
             \mu=\{(a,x),(b,y)\},\qquad
             \nu=\{(a,y),(b,x)\}
\]
are stable. The proposer-optimal one
is \(\mu\), whose stipulated utility
sum is 2; the other has sum 20.
Neither Pareto-dominates the other
for all four participants.
Thus stability and proposer optimality
do not imply maximum aggregate utility.
\end{hbtheorem}

The existing preference representation
in \cref{sec:p17-finite-preferences}
supplies ordinal comparison vocabulary.
It does not supply transferable money
units or interpersonal welfare weights.
The planned strategic-reasoning treatment
must also distinguish a Nash equilibrium,
which compares unilateral strategy changes
at a strategy profile, from an NM stable
set of imputations and from a stable
set of matching pairs.

\begin{hbdefinition}{def:p18-superadditive-cover}{The exact superadditive-cover construction}
For any real characteristic function \(v\)
on a finite \(N\), with \(v(\varnothing)=0\),
define
\[
 \widehat v(S)=\max_{\pi\text{ a partition of }S}
                              \sum_{K\in\pi}v(K).
\]
For the empty set the sole partition
is empty and the sum is zero.
The maximum exists because the
set of partitions is finite.
\end{hbdefinition}

\begin{hbtheorem}{thm:p18-superadditive-cover}{proved=annex:p18.superadditive.cover}{Preserving the complete dominance graph}
\(\widehat v\) is superadditive, satisfies
\(\widehat v(S)\geq v(S)\), and preserves
every singleton value. If
\(\widehat v(N)=v(N)\), the imputation
sets and their domination relations
for \(v\) and \(\widehat v\) are
identical. Consequently every argument
extension family and every NM solution
is unchanged. This gives an exact
bridge when a source initially presents
a characteristic function without
superadditivity.
\end{hbtheorem}

\begin{hbdefinition}{def:p18-lucas-game}{The exact ten-player characteristic function}
Let \(N=\{1,\ldots,10\}\),
\(O=\{1,3,5,7,9\}\) and
\(Q_j=\{2j-1,2j\}\), \(1\leq j\leq5\).
Lucas's original function \(\ell\) is
defined by the following complete list:
\[
\begin{aligned}
 \ell(N)&=5,&\ell(O)&=4,&\ell(Q_j)&=1,\\
 \ell(O\setminus\{i\})&=3
       &&(i\in\{1,3,5\}),\\
 \ell(O\setminus\{i,p\})&=2
       &&(i\in\{1,3,5\},\ p\in\{7,9\}),\\
 \ell(\{1,4,7,9\})&=\ell(\{3,6,7,9\})
                    =\ell(\{5,2,7,9\})=2.
\end{aligned}
\]
All other coalition values are zero.
Put
\[
 A_\ell=\{x\in\mathbb R^{10}:x_i\geq0,\ x(N)=5\},
 \qquad B_\ell=\{x\in A_\ell:x(Q_j)=1\text{ for all }j\}.
\]
Let \(c^0=(1,0,1,0,1,0,1,0,1,0)\);
for \(1\leq j\leq5\), let \(c^{2j}\)
replace the \(j\)-th pair \((1,0)\)
of \(c^0\) by \((0,1)\)
\cite[sec.~3, pp.~4--5]{Lucas1968NoSolutionProof}.
\end{hbdefinition}

\begin{hbtheorem}{thm:p18-lucas-core-foundation}{proved=annex:p18.lucas.core.foundation}{The complete source-game and core foundation}
\(\ell\) is not superadditive, but its
superadditive cover has the same
imputation set \(A_\ell\) and the
same domination graph. Their common
undominated core is
\[
 {\cal C}_\ell
   =\{x\in B_\ell:x(O)\geq4\}
   =\operatorname{conv}\{c^0,c^2,c^4,c^6,c^8,c^{10}\}.
\]
Thus the source's exact game can be
placed in the declared superadditive
domain without changing its NM
solutions. The full no-solution proof and its positive-coalition
preservation result are supplied in
\cref{thm:p18-lucas-no-solution,thm:p18-lucas-positive-cover}.
\end{hbtheorem}

The complete regional exclusion argument is supplied in
\cref{sec:p18-lucas-no-solution}. It rules out an NM solution
for both the original ten-player function and its superadditive cover.
